\section{Rows as homomorphic copies and partial gluings}\label{sec:rows}

\subsection{Every named row is a homomorphic carrier image}

The basic structural observation is immediate but decisive.

\begin{proposition}[Row-homomorphism principle]\label{prop:rowhom}
For each named $b\in B$, the map
\[
j_b:A\to \mathbb I_E^{\A},
\qquad
j_b(a)=[\alpha(b,a)]
\]
is a homomorphism of the native carrier signature.  It factors through
\[
\bar j_b:Q_E\to \mathbb I_E^{\A}.
\]
\end{proposition}

\begin{proof}
For a native $k$-ary $f$ and named $a_1,\ldots,a_k$, compatibility \eqref{eq:compat} and the row of $D_A$ give
\[
j_b(f^A(\bar a))
=
[f(\alpha(b,a_1),\ldots,\alpha(b,a_k))]
=
f^{\mathbb I_E}(j_b(a_1),\ldots,j_b(a_k)).
\]
If $a\equiv_{\theta_E^A}a'$, congruence under $\alpha(b,-)$ gives $j_b(a)=j_b(a')$, so $j_b$ factors through $Q_E$.
\end{proof}

Thus missing cells do not generate unrelated points.  For each operator $b$ they lie in a homomorphic image of the whole carrier quotient.  A supplied row $\alpha(b,a)=c$ identifies the row point $\bar j_b([a])$ with the distinguished carrier point $[c]$.  This is the local gluing picture that will be made exact by pushouts.

\subsection{Strong amalgamation in the unrestricted carrier signature}

Let $\Sigma_A^+$ denote the positive-arity native carrier signature, with the designated names omitted.  We use the category
\[
\Alg_{\Sigma_A^+}
\]
of all algebras of this fixed signature, not the variety generated by $A$.  This distinction matters: the ground diagram specifies the operation table of the named copy but does not impose every identity true in $A$ on arbitrary elements of a model.

We use the following standard free-amalgamation fact.

\begin{lemma}[Strong amalgamation]\label{lem:sap}
The category $\Alg_{\Sigma}$ of all algebras of a fixed finitary positive-arity signature has the strong amalgamation property.  If embeddings $S\hookrightarrow C$ and $S\hookrightarrow D$ are given, their pushout can be chosen so that both $C$ and $D$ embed and their images intersect exactly in the common image of $S$.
\end{lemma}


\begin{proof}
Identify the two copies of $S$ and put $U=C\sqcup_S D$. Construct normal
forms starting with the elements of $U$. For a native operation symbol $f$
and already constructed normal forms $x_1,\ldots,x_k$, evaluate the tuple
inside $C$ if all $x_i$ belong to $C$, or inside $D$ if all belong to $D$.
When both rules apply, all arguments lie in $S$, so the two answers agree.
When neither rule applies, retain $f(x_1,\ldots,x_k)$ as a new formal node.
Finite recursion gives a set of normal forms closed under these operations.

The inclusions of $C$ and $D$ are homomorphisms, remain injective, and have
intersection exactly $S$, since no two distinct base elements were
identified. Any two homomorphisms out of $C$ and $D$ agreeing on $S$ extend
uniquely by recursive evaluation of formal nodes. This is the pushout
universal property. The construction is the standard free amalgam for an
unrestricted signature; see also~\cite{Jezek2008,Zangurashvili2003}.
\end{proof}


\begin{remark}
Lemma~\ref{lem:sap} would be false in this form for an arbitrary nontrivial variety.  In a variety, additional identities can create dominion or amalgamation phenomena.  The clean kernel theory below is tied to the ground-diagram semantics of the present paper.
\end{remark}

\subsection{The raw row pushout}

Fix a congruence $\theta\in\Con(A)$ and put $Q=A/\theta$.  Fix a named $b\in B$, and let $R_b$ be the finite set of supplied row occurrences for $b$.  If $R_b\neq\varnothing$, let
\[
F_b=F_{\Sigma_A^+}(\{x_r:r\in R_b\})
\]
be the absolutely free carrier algebra on one generator for each occurrence.  Let $Q_b\cong Q$ be a formal row copy and $Q_*\cong Q$ the distinguished carrier copy.  If occurrence $r$ is
\[
\alpha(b,a_r)=c_r,
\]
define
\[
u_b:F_b\to Q_b,
\qquad u_b(x_r)=[a_r],
\]
and
\[
v_b:F_b\to Q_*,
\qquad v_b(x_r)=[c_r].
\]
Form the pushout
\begin{equation}\label{eq:rawpo}
\begin{array}{ccc}
F_b & \xrightarrow{u_b} & Q_b\\
{\scriptstyle v_b}\downarrow && \downarrow{\scriptstyle p_b}\\
Q_* & \xrightarrow{m_b} & P_b.
\end{array}
\end{equation}
If $R_b=\varnothing$, use the coproduct of the two disjoint copies of $Q$.
The same normal-form construction applies without a common subalgebra; the
positive-arity convention is relevant here. The row is automatically
protected at this stage.

Let
\[
S_b=\im u_b
=\langle[a_r]:r\in R_b\rangle_Q.
\]
The next theorem gives the exact local obstruction.

\begin{theorem}[Raw row protection criterion]\label{thm:raw}
Assume $R_b\neq\varnothing$.  The carrier leg $m_b:Q_*\to P_b$ in \eqref{eq:rawpo} is injective if and only if both of the following hold.
\begin{enumerate}[label=(\alph*)]
\item \emph{Homomorphic coherence:}
\[
\ker u_b\subseteq\ker v_b.
\]
Equivalently, the supplied values induce a unique homomorphism
\[
h_b:S_b\to Q
\]
with $h_bu_b=v_b$ after corestricting $u_b$ to $S_b$.
\item \emph{Kernel-extension stability:} with
\[
K_b=\ker h_b,
\qquad
\kappa_b=\Cg_Q(K_b),
\]
one has
\begin{equation}\label{eq:cepcondition}
\kappa_b\cap S_b^2=K_b.
\end{equation}
\end{enumerate}
\end{theorem}

\begin{proof}
Suppose first that $m_b$ is injective.  If $u_b(x)=u_b(y)$, then in the pushout
\[
m_bv_b(x)=p_bu_b(x)=p_bu_b(y)=m_bv_b(y).
\]
Injectivity gives $v_b(x)=v_b(y)$, so $\ker u_b\subseteq\ker v_b$.  Hence $v_b$ factors uniquely through $S_b$, defining $h_b$.

Write $u_b=i_be_b$ with $e_b:F_b\twoheadrightarrow S_b$ and $i_b:S_b\hookrightarrow Q_b$.  Since $e_b$ is surjective, a cocone over $(u_b,v_b)$ is the same thing as a cocone over $(i_b,h_b)$; thus \eqref{eq:rawpo} is also the pushout of
\[
S_b\hookrightarrow Q_b,
\qquad
h_b:S_b\to Q_*.
\]
Factor $h_b$ through its kernel:
\[
S_b\twoheadrightarrow S_b/K_b\hookrightarrow Q_*.
\]
Pushing the first map through $S_b\hookrightarrow Q_b$ quotients the row copy by
\[
\kappa_b=\Cg_Q(K_b).
\]
The induced map $S_b/K_b\to Q_b/\kappa_b$ is injective exactly when \eqref{eq:cepcondition} holds.  If it is injective, the remaining pushout is a strong amalgam of two embeddings, so Lemma~\ref{lem:sap} preserves the central carrier.  Conversely, if \eqref{eq:cepcondition} fails, two distinct $K_b$-classes from $S_b$ already merge in $Q_b/\kappa_b$; the pushout relation then forces their $h_b$-images to merge in $Q_*$, contradicting injectivity of $m_b$.
\end{proof}

The first obstruction includes ordinary value conflicts and all compatibility consequences on the subalgebra generated by the supplied inputs.  The second is a particular congruence-extension obstruction: even a genuine partial homomorphism can force additional equality if its kernel does not extend cleanly to the whole carrier.
